\section{Related Works} \label{sec:rel}

The foundation of our work is the VQ-VAE framework of \cite{VQVAE}. Our prior network is based on Gated PixelCNN \cite{pixelcnn} augmented with self-attention \cite{Vaswani2017}, as proposed in \cite{PixelSnail}. 
\par
BigGAN \cite{biggan} is currently state-of-the-art in FID and Inception scores, and produces high quality high-resolution images. The improvements in BigGAN come mostly from incorporating architectural advances such as self-attention, better stabilization methods, scaling up the model on TPUs and a mechanism to trade-off sample diversity with sample quality. In our work we also investigate how the addition of some of these elements, in particular self-attention and compute scale, indeed also improve the quality of samples of VQ-VAE models.
\par
Recent work has also been proposed to generate high resolution images with likelihood based models include Subscale Pixel Networks of \cite{SPN}. Similar to the parallel multi-scale model introduced in \cite{reed2017parallel}, SPN imposes a partitioning on the spatial dimensions, but unlike \cite{reed2017parallel}, SPN does not make the corresponding independence assumptions, whereby it trades sampling speed with density estimation performance and sample quality.
\par
Hierarchical latent variables have been proposed in e.g. \cite{Rezende2014}. Specifically for VQ-VAE, \cite{Dieleman2018} uses a hierarchy of latent codes for modeling and
generating music using a WaveNet decoder. The specifics of the encoding is however different from ours: in our work, the bottom levels of hierarchy do not exclusively refine the information encoded by the top level, but they extract complementary information at each level, as discussed in \sref{sec:vq-ms}. Additionally, as we are using simple, feed-forward decoders and optimizing mean squared error in the pixels, our model does not suffer from, and thus needs no mitigation for, the hierarchy collapse problems detailed in \cite{Dieleman2018}. Concurrent to our work, \cite{defauw2019} extends \cite{Dieleman2018} for generating high-resolution images. The primary difference to our work is the use of autoregressive decoders in the pixel space. In contrast, for reasons detailed in \sref{sec:method}, we use autoregressive models exclusively as priors in the compressed latent space, which simplifies the model and greatly improves sampling speed. Additionally, the same differences with \cite{Dieleman2018} outlined above also exist between our method and \cite{defauw2019}.
\par
Improving sample quality by rejection sampling has been previously explored for GANs \cite{azadi2018discriminator} as well as for VAEs \cite{Bauer18} which combines a learned rejecting sampling proposal with the prior in order to reduce its gap with the aggregate posterior.